% Shaurya's section, transcribed verbatim from his compiled report (report-1.pdf).
\section{Tuning for This Machine}
Table~\ref{tab:rules} lists the rules; the ones that need explanation follow.

\textbf{Vector size depends on the work, not on one cache level.} A vector's buffers live in
L1, L2 or L3 depending on $B$. Work per value favours small vectors and a fixed cost per vector
favours large ones, and a synthetic benchmark shows either can win: random reads inside the
vector are fastest with L1-sized buffers (1.34 vs 1.69 and 2.55\,ns per value for L2 and L3
sizes), the same with 2\,$\mu$s per vector with L2-sized ones (1.90\,ns), and light work with
20\,$\mu$s per vector with L3-sized ones (1.69 vs 12.05\,ns for L1). Our scans read their buffers
sequentially, and an L2 hit streams nearly as fast as an L1 hit (130 vs 122\,GB/s), so L1-sized
vectors cut L1D misses but not time, while small vectors slow decoding. The rule keeps every byte
a vector touches in a quarter of L2.

\begin{table*}[t]
\caption{Each machine parameter decides one code choice. Knobs are computed from sysfs at
start-up and are off in the untuned baseline. Times are from DRAM unless noted.}
\label{tab:rules}
\small
\begin{tabularx}{\textwidth}{>{\raggedright\arraybackslash}p{0.17\textwidth}>{\raggedright\arraybackslash}p{0.12\textwidth}Xp{0.21\textwidth}}
\toprule
Parameter & Decides & Rule & Effect \\
\midrule
Line size, value width; where the table lives & early or late & late while $f(s) <
\varphi(1-\sigma/g)$; one constant pair for DRAM, one for tables that fit in L3 & within 1\% of
the better fixed strategy on DSM; \code{SELECT *} in L3 worst case 2.6$\times$ $\to$
1.3$\times$ \\
Cache sizes vs work per value & vector size & buffers in $\frac{1}{4}$ of L2 for our scans &
decoding SUM 11.9 $\to$ 9.9\,ms vs an L1-sized vector \\
L1: 64 sets $\times$ 8 ways, 4\,KB stride & DSM column starts, PAX minipage stride & column $c$
starts $17c$ lines in; minipages of an odd number of lines & early \code{SELECT *} at 1\%: DSM
62.6 $\to$ 22.9\,ms, PAX (64\,KB) 65.6 $\to$ 27.8\,ms \\
Write-allocate (12 vs 40\,GB/s with streaming stores) & how output is written & normal stores
into a cache-resident stage, copied out with streaming stores & \code{SELECT *} of every row
1.2--1.9$\times$ (DSM early 71.2 $\to$ 37.9\,ms) \\
Prefetcher stops at 4\,KB pages & PAX page size & 256\,KB pages & PAX SUM(c0) 4.5 $\to$
1.2\,ms \\
dTLB reach (2048 entries $\times$ 4\,KB = 8\,MB) & OS page size & 2\,MB pages for buffers
$\geq$ 2\,MB & late SUM of 16 columns 23.3 $\to$ 18.7\,ms \\
AVX2, no compress instruction & kernels, code width & SIMD selection vector, fused
unpack-and-add, byte-aligned dictionary codes & see Sections 5 and 6 \\
\bottomrule
\end{tabularx}
\end{table*}

\textbf{Cache sets.} Streams read in lockstep must not share an L1 set. Staggering DSM columns by
17 lines (coprime to the 64 sets) puts up to 64 columns in distinct sets; shortening PAX
minipages to an odd number of lines (64 $\to$ 63 at 64\,KB, 1.6\% of space) does the same for
large pages. Padding by 64 bytes does not, because the stride stays $P/16$. The gain is large
only where early reads 16 streams and little is written: at 100\% or for SUM of 16 attributes,
the misses halve but time improves only 1.1$\times$.

\textbf{Output writes.} Streaming stores skip the read before write, but only for whole lines
written in order, and our writes are scattered (early overwrites a failing row's slot, late
writes one attribute at a time). So each vector writes into a cache-resident stage with normal
stores, which absorb the scattered writes, and the stage is copied out with streaming stores. A
stage copied with \code{memcpy} is up to 36\% \emph{slower}: the gain is the streaming copy.

\textbf{Huge pages.} With 2\,MB pages, dTLB misses for SUM of 16 DSM columns fall from 0.031 to
0.008 per row for both strategies, but only late gets faster (20\% vs 3\%): early's 16
interleaved columns keep loads going during a page walk, while late's single stream waits.

\textbf{Together}, over six queries $\times$ three layouts from DRAM, against an untuned engine
given the best of four fixed strategies per cell, the tuned engine is \textbf{1.46$\times$}
faster (geometric mean; NSM 1.11$\times$, DSM 1.41$\times$, PAX 1.96$\times$), up to 3.8$\times$
for PAX SUM(c0). \code{SELECT *} of every row gains 1.7--2.5$\times$, after which NSM no longer
wins it (36.9\,ms, DSM 34.6, PAX 31.7). Queries that already move few lines gain little (DSM
SUM(c0): 1.03$\times$).

\textbf{Tried and dropped.} Instead of a threshold, auto can time both strategies per span and
run the cheaper one, with \code{-{}-auto-measure}. This bounds the worst case in L3 to
1.16$\times$, but costs 2--5\% where the threshold is right and 5--15\% in cache; the tuned run
above used it. \emph{One gather pass for all columns} took 0.97--1.11 times as long below 5\%,
because the per-column loop already keeps many misses in flight. \emph{SUM over dictionary codes
by counting} was 10--25\% slower than one lookup per row, because counter updates to the same
code wait on each other through memory.
